% Full elementary-interval and fusion proofs for Kamp III.2, Lemma 4.

\section{Kamp's conjunction lemma for interval conditions}
\label{sec:annex-p09-kamp-iii2-lemma4}

\begin{annexitem}{annex:p9.kamp.elementary.intervals}{thm:p09-kamp-elementary-intervals}{Elementary labels and their finite substitution presentations}
The source classes and remarks R5--R6 occur on printed pp.~67--68.
The finite-chain interpretation used here is the one proved in
\cref{annex:p9.kamp.iii2.lemma1}; proper substitution has the semantics
proved in \cref{annex:p9.kamp.proper.substitution}
\cite[ch.~III, pp.~63--68]{Kamp1968}.

\begin{proof}
Fix the equality-free language \(L\). All renamings of bound variables
below are alphabetic changes preserving their scopes. In particular, a
free endpoint is never replaced at an occurrence bound by an internal
quantifier.

\paragraph{Changing the name of the sole free parameter.}
A finite construction in \(\operatorname{E}(v)\) gives one in
\(\operatorname{E}(w)\) by renaming the free parameter \(v\) to
\(w\), after making its bound names fresh. For an
\(\operatorname{At}\) seed this changes the argument of each literal,
or of the two occurrences in its particular \(P_0\) tautology.
The resulting seed belongs to \(\operatorname{At}(w)\).
For a basic past seed, its free right endpoint changes from \(v\) to
\(w\); its left endpoint and internal variables remain bound and can
first be made fresh. The result is the same basic-past constructor at
\(w\). For a basic future seed it is the free left endpoint that
changes. For a limiting past seed the renamed parameter is the fixed
right boundary in every outer cut and interior-witness guard; all cut
and witness variables remain bound. The result is precisely the
limiting-past constructor at \(w\). The limiting future constructor
has the fixed left boundary and has the same renaming property. These
arguments include the one-predicate limit seeds.

Negation, conjunction and disjunction commute with the free-parameter
renaming. At a proper-substitution node, first make the bound variables
of both children fresh from \(v,w\) and from their insertion contexts.
At every original target atom, both orders of construction insert the
same replacement with the same free argument; only selected bound names
can differ. The alphabetic-variation result in the preceding substitution
proof identifies those outputs. Induction on the finite construction
therefore proves the parameter-renaming assertion for every
\(\operatorname{E}\) member. No substitution for a bound occurrence
of the parameter was used.

\paragraph{Following a substitution history at every original slot.}
Start with a basic pattern and its finite occurrence word. At an original
occurrence of a unary predicate \(Q\), attach the label \(Q(w)\).
It belongs to \(\operatorname{At}(w)\), including when the same
predicate occurs at several slots. The source pattern is exactly the
appropriate recursion \(\mathcal B_p\) or \(\mathcal B_f\) with
these initial labels, up to its bound names.

Suppose a history step properly substitutes \(\psi\in
\operatorname{E}(v)\) for \(P_i\). Replace its parameter by \(w\)
using the preceding paragraph, and apply the proper substitution for
\(P_i\) to every current label \(\theta\in\operatorname{E}(w)\).
Each new label is in \(\operatorname{E}(w)\) by that class's actual
substitution constructor. Proper substitution leaves every order atom
and every outer Boolean and quantifier constructor in the basic skeleton
unchanged. Its only insertions are inside the current label formulas.
Freshness gives the same inserted formula at each target atom whether
the label is transformed first or at its location in the pattern.
Hence the transformed pattern is the same labelled recursion with the
new labels, up to alphabetic variation.

Induct over the entire finite substitution history. This proves that
every elementary interval formula has the asserted elementary labelled
word. A predicate inserted by an earlier step is part of the current
label and is visited by a later step if it is then the target. A
predicate inside the replacement of the current step is not visited
again in that step. Thus the construction also handles a replacement
containing its own target, shared letters, absent targets and repeated
occurrences. The argument respects the order of the history rather than
treating it as simultaneous substitution.

\paragraph{Presenting any labelled word by the source's initial-segment list.}
Take a word \(W\) with \(2m+1\) elementary labels. Only finitely many
predicate letters occur in all these finite formulas. Choose distinct
placeholder letters \(Q_1,\ldots,Q_{2m+1}\), all with indices larger
than every predicate index in the labels, and use them in order in a
basic past pattern. Such a pattern exists by the explicit word
construction in \cref{annex:p9.kamp.iii2.lemma1}. Each placeholder
occurs exactly once.

In increasing predicate-index order, substitute the corresponding label
properly for every placeholder. The labels contain no placeholder, so
a later substitution cannot change an earlier inserted label. For every
unused predicate index up through the last placeholder, use the identity
replacement \(P_i(w)\). Those replacements belong to
\(\operatorname{E}(w)\) and change no formula. The resulting full
list has exactly the source form with targets \(P_0,\ldots,P_N\).
Its output is \(\mathcal B_p(W;x,y)\), up to fresh bound names.
This exhibits an \(\operatorname{Ep}(x,y)\) presentation without
appealing to closure under semantic equivalence.

For the future pattern, choose the basic future pattern on the same
placeholder word. Substitute the same labels in the same order. Its
output is exactly the recursion \(\mathcal B_f(W;x,y)\), up to
bound names, so it has an \(\operatorname{Ef}(x,y)\) presentation.

Every initial-segment history is a finite history. Conversely the
slot argument transforms any finite history into a labelled word, and
the placeholder argument presents that word by an initial-segment
history. This proves the equivalence of the two descriptions of the
elementary interval classes used in the body. It also justifies appending
a substitution for an arbitrary target: append it to a finite history,
then use this conversion if an initial-segment presentation is required.
This is R6 for both classes.

\paragraph{The realization equivalence.}
Fix any nonempty linear structure \(\mathcal M\) and ordered endpoints
\(a<b\). Interpret each fresh placeholder \(Q_s\) by the truth set
of its label \(\theta_s(w)\) in \(\mathcal M\). A label has exactly
one free variable, so this set is independent of the ambient assignment
to other variables, by the variable-dependence proof in
\cref{annex:p9.kamp.proper.substitution}. The underlying order is unchanged.

The full first-order substitution identity in that annex shows that the
substituted pattern is true in \(\mathcal M\) exactly when the basic
placeholder pattern is true in this modified unary interpretation.
Since no label contains a placeholder, the labels are evaluated in
\(\mathcal M\) itself and the successive substitutions have precisely
these fixed truth sets. The basic-pattern realization equivalence in
\cref{annex:p9.kamp.iii2.lemma1} now says that its endpoints can be
joined by a strictly increasing sequence satisfying every point label
and every open-interval label, exactly as asserted in the body.

Concretely, at one interval the statement is the two endpoint formulas
and the single universal gap condition. At a past constructor the new
existential witness supplies the next point after the left endpoint,
its new universal guard supplies the first gap, and the remaining
realization begins at that witness. Conversely, the second point of a
realization supplies this witness and the remaining suffix. At a future
constructor the existential supplies the penultimate point, the new
universal guard supplies the last gap, and the earlier realization ends
there. Conversely, the penultimate point of a realization supplies that
witness and the preceding prefix. These are exactly the two basic
constructor cases of the cited induction. Interpreting placeholders by
arbitrary label truth sets therefore preserves every case, even when a
label's own quantifiers inspect points outside the displayed interval.

Proper substitution preserves the free-variable set of the outer
formula. The basic patterns have exactly \(x,y\) free, so both
elementary classes and both labelled recursions have those endpoints
free. The realization assertion has required \(a<b\) throughout;
it has not added an order conjunct to the source grammar.

\paragraph{The guarded one-variable closures.}
For \(\phi\in\operatorname{Ep}(x,y)\), start with its basic
past skeleton \(\phi_0\). Its guarded closure
\(\exists x\,(x<y\land\phi_0)\) belongs to
\(\operatorname{Bp}(y)\), hence to \(\operatorname{E}(y)\).
Apply to it the finite substitution history of \(\phi\), first
renaming every replacement's sole free parameter to \(y\).
Each step is an allowed \(\operatorname{E}(y)\) substitution.
The guard has only order atoms, so it is unchanged. The outer existential
still binds the free argument occurrences at the left endpoint, and
freshness protects all internal bound variables. Thus the resulting
formula is \(\exists x\,(x<y\land\phi)\), up to the allowed bound
renaming. This proves the past half of R5 by an actual finite construction.

For \(\phi\in\operatorname{Ef}(x,y)\), start instead with the
basic future skeleton \(\phi_0\). Its guarded closure
\(\exists y\,(x<y\land\phi_0)\) is in
\(\operatorname{Bf}(x)\), hence in \(\operatorname{E}(x)\).
Rename replacement parameters to \(x\) and follow its history.
The order guard remains unchanged, the outer quantifier binds precisely
the right-endpoint arguments, and inserted bound variables are fresh.
The result is the required guarded future formula in
\(\operatorname{E}(x)\), up to bound names. This proves the other
half of R5. Missing past or future endpoints make the respective
existentials false; empty internal open intervals retain their ordinary
vacuous universal truth. All preceding arguments hold over
\(\mathsf{Lin}\), without completeness.
\end{proof}
\end{annexitem}

\begin{annexitem}{annex:p9.kamp.interval.substitution}{thm:p09-kamp-interval-substitution}{The substitution and endpoint remarks R7--R8}
These two remarks are printed immediately after the definition of
\(I(x,y)\) \cite[ch.~III, p.~68]{Kamp1968}.

\begin{proof}
Let \(\phi\) have a representation
\(\mathsf{Com}\models x<y\to(\phi\leftrightarrow\bigvee_i\theta_i)\),
with \(\theta_i\in\operatorname{Ep}(x,y)\), and let
\(\psi\in\operatorname{E}(v)\).

Fix a structure \(\mathcal M\) over a complete linear order.
Interpret \(P_j\) instead by the set of points satisfying
\(\psi(v)\) in \(\mathcal M\); leave the domain, order and all other
unary interpretations unchanged. The resulting structure is still over
the same complete order, so the given \(\mathsf{Com}\) equivalence
applies to it. The first-order substitution identity proved in
\cref{annex:p9.kamp.proper.substitution} transfers that equivalence back
to \(\mathcal M\) as
\[
 x<y\longrightarrow
 \left(\operatorname{psub}_{P_j}(\phi,\psi)
   \longleftrightarrow
        \bigvee_i\operatorname{psub}_{P_j}(\theta_i,\psi)\right).
\]
The order guard is unchanged because its binary predicate is not a
substitution target. R6, just proved, places each resulting disjunct in
\(\operatorname{Ep}(x,y)\). Proper substitution preserves all free
variables, so the new left side has no extra free parameter. This proves
R7. The construction evaluates \(\psi\) in the original structure
even if it contains \(P_j\); hence no fixed-point interpretation is
involved.

For R8, take \(\theta\in\operatorname{E}(x)\), and replace its
sole free parameter by \(w\), obtaining a label \(\theta_w\).
The first annex item proves this parameter change preserves the
elementary class. Use the two words
\[
        (\theta_w,T^K_w,T^K_w),
                 \qquad (T^K_w,T^K_w,\theta_w).
\]
The first past pattern, at endpoints \((x,y)\), is
\[
 \theta(x)\land\forall z\,(x<z<y\to T^K_z)\land T^K_y.
\]
Its interval and last-point factors are tautologies, so it agrees with
\(\theta(x)\) at every ordered pair and has an
\(\operatorname{Ep}(x,y)\) presentation. This gives \(I(x,y)\).
For the second word put the endpoints in the order \((y,x)\). Its
formula is
\[
 T^K_y\land\forall z\,(y<z<x\to T^K_z)\land\theta(x).
\]
It agrees with \(\theta(x)\) whenever \(y<x\), and has an
\(\operatorname{Ep}(y,x)\) presentation. This gives \(I(y,x)\).
These two displayed equivalences already hold on every linear order;
restrict them to \(\mathsf{Com}\) for the stated properties.
The truth labels are the source's particular formulas containing
\(P_0\), not undeclared constants or predicates.
\end{proof}
\end{annexitem}

\begin{annexitem}{annex:p9.kamp.iii2.lemma4}{thm:p09-kamp-iii2-lemma4}{Every finite fusion and both directions of the conjunction construction}
The source proof first distributes the given interval representations,
then passes to distinct predicate occurrences, and finally fuses their
endpoint sequences. We prove the fusion identity with elementary labels
and give the precise connection to that source construction
\cite[ch.~III, pp.~73--78]{Kamp1968}.

\begin{proof}
Fix words \(W,V\) with respectively \(m,n\geq1\) intervals, and
use the tags, fusions and slot labels defined in the body. Their common
label parameter is \(w\); all binding names in the interval formulas
are fresh. We first work over an arbitrary nonempty strict linear order.
In the semantic arguments, \(\theta(t)\) means that the label holds
with its free parameter assigned the point \(t\). Likewise,
\(\chi_F(a,b)\) and \(\mathcal B_p(W;a,b)\) mean satisfaction of their
formulas with \(x\) assigned \(a\) and \(y\) assigned \(b\).

\paragraph{The finite fusion data are well defined.}
For a total preorder, mutual relatedness is an equivalence relation:
reflexivity and symmetry follow directly from its definition; if
\(a\sim_F b\) and \(b\sim_F c\), transitivity in both directions
gives \(a\sim_F c\). Its equivalence classes inherit a linear order.
Indeed comparison comes from totality; replacing representatives preserves
it by transitivity; and two distinct classes cannot be related in both
directions. Thus the finite class list \(C_1,\ldots,C_{r_F}\) and
the map \(h_F\) are well defined.

The restriction of \(F\) to either input is its strict sequence
order together with reflexive pairs. Consequently \(h_F(e_i)\)
strictly increases with \(i\), and \(h_F(f_j)\) strictly increases
with \(j\). A class contains at most one point tag from each input.
The two first tags form the first class and the two last tags form the
last class, by the endpoint identifications and the order of every
other tag within its own sequence. Since each input has at least two
points, these classes differ and \(r_F\geq2\).

There are finitely many binary relations on the finite tag set. Enumerate
them and retain precisely the reflexive, transitive, total relations
with the required restrictions and endpoint identifications. This gives
a finite enumeration of exactly \(\mathcal F_{m,n}\), using tests on
finite sets. At least one such fusion exists: identify the two first
tags and the two last tags, place every interior \(e\)-tag before
every interior \(f\)-tag, and preserve each input's order.

The source's representative convention gives the same length. Each
class either contains an \(e\)-tag, in which case that tag is its
unique representative from the first sequence, or contains only one
\(f\)-tag, which is independent in the source's sense. Conversely
every such selected tag lies in one class and no two selected tags
lie in the same class. Hence the length counts the classes, not the
number of original tags or the number of alternating slots.

\paragraph{Every slot has a label and every constructed formula is elementary.}
An odd slot \(2k-1\) represents the nonempty class \(C_k\).
Choose any point tag in this class. Its required-slot singleton is
\(\{2k-1\}\), so this slot has at least its point label.

For an even slot \(2k\), where \(1\leq k<r_F\), the increasing
list
\(h_F(e_0)=1<\cdots<h_F(e_m)=r_F\) has a unique consecutive
pair with
\(h_F(e_{i-1})\leq k<h_F(e_i)\).
For example choose the greatest index whose position is at most \(k\);
it is less than \(m\), and the next index gives the pair.
It follows that
\[
             2h_F(e_{i-1})-1<2k<2h_F(e_i)-1,
\]
so \(2k\in f_F(b_i)\). The increasing \(f\)-position list
gives a unique \(d_j\) in the same way. Thus every even slot has
an interval label from each input, and every slot's conjunction is
nonempty. A finite nonempty conjunction of \(\operatorname{E}(w)\)
formulas belongs to \(\operatorname{E}(w)\), by iterating its binary
conjunction constructor. Hence every \(\rho_{F,s}\) is an elementary
label. The first annex item gives an actual
\(\operatorname{Ep}(x,y)\) presentation for \(\chi_F\).

An interval tag's coverage formula excludes its own endpoint slots,
includes every fused point strictly between those endpoints, and includes
every fused gap between them. This describes required positive truths.
It does not constrain the truth of that label at slots outside this set.

\paragraph{A fused realization yields the first input realization.}
Fix endpoints \(a<b\) in a linear structure \(\mathcal M\), and
suppose \(\chi_F(a,b)\) holds. By the labelled realization equivalence,
choose
\[
                      a=u_1<\cdots<u_{r_F}=b
\]
realizing all its slot labels. For the first input put
\(a_i=u_{h_F(e_i)}\), for \(0\leq i\leq m\).
The indices strictly increase, and the first and last are 1 and
\(r_F\), so this is a strictly increasing sequence from \(a\)
to \(b\). At its \(i\)-th point the odd slot
\(2h_F(e_i)-1\) contains \(\alpha_i\) as a conjunct. Thus
\(\alpha_i(a_i)\) holds.

It remains to check the whole original interval for \(\beta_i\).
Write \(p=h_F(e_{i-1})\), \(q=h_F(e_i)\), so \(p<q\), and
take any \(t\) with \(u_p<t<u_q\). There are exactly two possible
kinds of position relative to this finite chain. Either \(t=u_k\)
for some \(p<k<q\), or
\(u_k<t<u_{k+1}\) for some \(p\leq k<q\). To see exhaustion
in the latter case, choose the greatest index with \(u_k<t\).
The index exists because \(u_p<t\), is smaller than \(q\), and
its maximality and the assumption that \(t\) is no chain point give
\(t<u_{k+1}\).

In the point case,
\(2p-1<2k-1<2q-1\), so \(\beta_i\) is a conjunct of
\(\rho_{F,2k-1}\). Truth of this point label gives \(\beta_i(t)\).
In the gap case,
\(2p-1<2k<2q-1\), so \(\beta_i\) is a conjunct of
\(\rho_{F,2k}\). The realization requires that conjunction at every
point of \((u_k,u_{k+1})\), hence at \(t\). Thus every point of
the original interval satisfies \(\beta_i\).
This proves every point and interval condition of \(W\), so
\(\mathcal B_p(W;a,b)\) holds.

\paragraph{A fused realization yields the second input realization.}
Put \(c_j=u_{h_F(f_j)}\), for \(0\leq j\leq n\). Its indices
strictly increase from 1 to \(r_F\). At each selected point the
singleton \(f_F(f_j)\) puts \(\gamma_j\) in its odd slot label,
so \(\gamma_j(c_j)\) holds.

For an interval of the second input, put
\(p=h_F(f_{j-1})<q=h_F(f_j)\), and take
\(u_p<t<u_q\). If \(t=u_k\) with \(p<k<q\), its odd slot
satisfies \(2p-1<2k-1<2q-1\), so \(\delta_j(t)\) follows
from its point conjunction. Otherwise the finite-chain argument in the
preceding paragraph gives \(u_k<t<u_{k+1}\) with \(p\leq k<q\).
Then \(2p-1<2k<2q-1\), so the gap conjunction contains
\(\delta_j\), and its universal truth gives \(\delta_j(t)\).
Every interval condition of \(V\) follows. The selected endpoints
are \(a,b\), so \(\mathcal B_p(V;a,b)\) holds. We have proved
that each individual fusion disjunct implies both input patterns.

\paragraph{Two input realizations determine an allowed fusion.}
Conversely suppose both input patterns hold at \(a<b\). Choose their
realizations
\[
 a=a_0<\cdots<a_m=b,
                  \qquad a=c_0<\cdots<c_n=b.
\]
Assign the point \(a_i\) to tag \(e_i\) and the point \(c_j\)
to tag \(f_j\). Define \(sFt\) exactly when the assigned point of
\(s\) is at most the assigned point of \(t\), where at most means
strictly earlier or identical. This relation is reflexive, transitive
and total. The strictness of each realization gives exactly the required
restriction on its own tags. The common endpoints give
\(e_0\sim_F f_0\) and \(e_m\sim_F f_n\). Thus this is one
of the enumerated fusions.

Its equivalence classes correspond exactly to equal assigned points.
List the distinct points of the union of the two realizations as
\(a=u_1<\cdots<u_{r_F}=b\). Then
\(a_i=u_{h_F(e_i)}\) and \(c_j=u_{h_F(f_j)}\).
This uses only already chosen witness points; it does not posit a
point in a missing interval or take an order-theoretic limit.

\paragraph{The induced realization satisfies every odd slot.}
Fix a point slot \(2k-1\), and consider any of its required tags.
If the tag is \(e_i\), its singleton says \(h_F(e_i)=k\), so
\(u_k=a_i\) and the first input realization supplies
\(\alpha_i(u_k)\). If it is \(f_j\), then \(u_k=c_j\) and
the second realization supplies \(\gamma_j(u_k)\).

If the tag is \(b_i\), the membership
\(2k-1\in f_F(b_i)\) says
\(h_F(e_{i-1})<k<h_F(e_i)\). Hence
\(a_{i-1}<u_k<a_i\), and the universal interval condition of
the first realization supplies \(\beta_i(u_k)\).
If the tag is \(d_j\), the same integer inequalities using the
second input give \(c_{j-1}<u_k<c_j\), so
\(\delta_j(u_k)\) follows from its interval condition.
These exhaust the tag types, proving every conjunct of
\(\rho_{F,2k-1}\) at \(u_k\).

\paragraph{The induced realization satisfies every even slot.}
Fix \(2k\) with \(1\leq k<r_F\), and take any
\(u_k<t<u_{k+1}\). A point tag's singleton is odd, so only interval
tags can require this slot. If \(b_i\) requires it, its defining
inequality is equivalent to
\(h_F(e_{i-1})\leq k<h_F(e_i)\). Therefore
\[
 a_{i-1}\leq u_k<t<u_{k+1}\leq a_i,
\]
which places \(t\) strictly in \((a_{i-1},a_i)\). The first
input supplies \(\beta_i(t)\). If \(d_j\) requires the slot,
the inequality is \(h_F(f_{j-1})\leq k<h_F(f_j)\), and
\(c_{j-1}\leq u_k<t<u_{k+1}\leq c_j\) similarly gives
\(\delta_j(t)\). Thus every required conjunct holds at every point
of the gap. This proves the even-slot universal condition; it also
covers an empty gap, for which there is no \(t\) to test.

Together the two slot verifications give a realization of the whole
fused word, so \(\chi_F(a,b)\) holds. This proves that the
conjunction of the inputs implies at least one enumerated disjunct.
Combined with the first direction it proves the complete
\(\mathsf{Lin}\) fusion identity in the body.

\paragraph{Connection with the source's distinct-predicate construction.}
Kamp first presents each elementary input as a basic pattern with
distinct predicate occurrences, followed by substitutions. The slot
history and fresh-placeholder argument in the first annex item provides
this presentation explicitly, choosing the two placeholder alphabets
disjoint and absent from all replacement labels. Its occurrence
sequence is the source's \(S(\phi')\). The subsequence of point
positions, which are the occurrences outside the universal gap
quantifiers, is the source's \(\operatorname{Se}(\phi')\).
For our first skeleton it consists of its placeholders at
\(e_0,\ldots,e_m\); for the second it consists of those at
\(f_0,\ldots,f_n\). These sequences form a pure pair even when
some final labels coincide.

For a source fusion of these skeleton sequences, identify its distinct
predicate occurrences with our tags. The endpoint requirements in source
display~(19) are exactly our two endpoint equivalences. The map called
\(f_i\) on p.~76 assigns an endpoint occurrence its odd slot and
an interval occurrence all slots strictly between the odd slots of its
two endpoints. These are precisely the four clauses of \(f_F\).
The source's conjunction at each slot consists of all the predicate
letters requiring it, so substituting the elementary labels into those
conjunctions gives exactly our \(\rho_{F,s}\).

Conversely, choose another fresh basic skeleton on the fused slot
word. It is in \(\operatorname{Bp}_{r_F-1}(x,y)\), since the word
has \(r_F\) point slots and \(r_F-1\) interval slots. Substituting
these finite conjunctions gives our \(\chi_F\).
For each substituted elementary input the corresponding finite history
has already been absorbed in its labels. Freshness of all skeleton
letters from all labels ensures that the remaining substitutions leave
earlier inserted labels unchanged. The first-order substitution identity
then transfers the pure-predicate fusion identity to the labelled one.
Thus the direct labelled proof supplies the same construction as the
source reduction and gives all of its omitted point and gap checks.
It does not assume any simultaneous or recursively self-applying
predicate replacement.

\paragraph{Using the representations on their stated class of orders.}
Now let \(\phi,\psi\) have property \(I(x,y)\), with finite
nonempty representations \(\bigvee_{i=1}^{M}\phi_i\) and
\(\bigvee_{j=1}^{N}\psi_j\). By definition, their guarded
equivalences hold over \(\mathsf{Com}\). Fix any structure over a
complete order and any \(a<b\). Boolean distribution gives
\[
 (\phi\land\psi)(a,b)
 \quad\Longleftrightarrow\quad
             \bigvee_{i=1}^{M}\bigvee_{j=1}^{N}
                         (\phi_i\land\psi_j)(a,b).
\]
For each pair, the first annex item obtains its two elementary labelled
words. The proved \(\mathsf{Lin}\) fusion identity holds on this
particular complete order because every complete linear order is a
linear order. Replace the pair by the finite disjunction of its fusion
formulas. Each resulting formula has an
\(\operatorname{Ep}(x,y)\) presentation. The disjunction over all
pairs and their finite fusion lists is finite. Its guarded equivalence
to \(\phi\land\psi\) on all complete orders is exactly property
\(I(x,y)\).

If an input representation was written as an empty disjunction, the
false elementary pattern in the body's definition replaces it by a
one-disjunct representation. Its first label is false at every point,
so no realization exists; the preceding argument applies unchanged.
No free variable other than possibly \(x,y\) enters the original
conjunction, and every produced interval disjunct has exactly its two
endpoints free. Assignments with \(x\not<y\) make the guard false,
so no assertion about their unguarded truth is needed.

This is the complete application of the source's Lemma~4. The opening
representation premises are used on their defining \(\mathsf{Com}\)
scope. Only the elementary fusion identity was proved on arbitrary
linear orders; nothing in this passage derives an all-linear-order
equivalence from one given solely on complete orders.

\paragraph{Boundary cases and the displayed examples.}
For \(m=n=1\), both inputs have only their two endpoints. The endpoint
identifications leave exactly two classes and one fusion. Its word is
\[
        (\alpha_0\land\gamma_0,\
           \beta_1\land\delta_1,\ \alpha_1\land\gamma_1).
\]
At two adjacent points the middle requirement is vacuous on both inputs
and on the fused word. Neither the absence of intermediate points nor
the presence of global endpoints changes this calculation. If one
input has further point positions, its strict inequalities require
that many actual distinct points; the fusion cannot collapse two of
its own positions or manufacture the required witnesses. A singleton
order has no ordered endpoint assignment and therefore satisfies every
guarded identity here vacuously.

For the two one-interior-point words in the body, a fusion must put the
two interior tags in strict order in one direction, identify them,
or put them in strict order in the other direction. There are no other
tags besides the already identified endpoints. Thus the three displayed
words exhaust the fusion list. If \(a<c\), the \(A_1\) point lies
inside the first \(D_1\) interval and the \(C_1\) point inside the
second \(B_2\) interval. The three gaps respectively require
\(B_1\land D_1\), \(B_2\land D_1\), and
\(B_2\land D_2\), giving exactly that row. If \(a=c\), the
middle point has both point labels; it is in neither of the adjacent
open intervals at their shared endpoint. The two gaps give the middle
row. If \(c<a\), the \(C_1\) point is in the first \(B_1\)
interval, the \(A_1\) point is in the second \(D_2\) interval,
and the three gaps give the last row. The endpoint point labels are
conjoined in every case.

On \(0<1<2\) with both middle labels true just at 1 and all other
labels true everywhere, each input has the realization \((0,1,2)\).
The coincident row has this same realization. Either strict-order row
would require two different interior points, which this interval lacks,
so neither is true. On \(0<1<2<3\) in the second example, the first
input must choose 1 for its middle point, and then requires the false
\(B_2\) at 2. It fails. The second input can choose 2 and satisfies
its conditions. In the \(a<c\) fusion the point label
\(C_1\land B_2\) at 2 is false. If that particular \(B_2\)
requirement were omitted, the sequence \((0,1,2,3)\) would satisfy
all the remaining point labels, while all its three open gaps would
be empty and hence would satisfy their gap conditions. This would be
a false positive. The odd slots inside an original interval are
therefore necessary, even on finite complete orders.

There is likewise no negative requirement on predicates outside their
assigned slots. On a two-point order with every input predicate true
everywhere, the fused one-interval word is true. An interval predicate
is nevertheless true at the endpoints too, although its required-slot
set contains only the middle even slot. Truth at a point therefore
does not imply membership of that point's slot in its required-slot
set. The valid proof uses required-slot membership to obtain truth,
as in the two forward realization paragraphs.

Finally, the input histories and words are finite, the fusion set is
finite, and each output slot is a finite conjunction of existing labels.
Each displayed recursion removes one interval, and every membership
presentation uses finitely many fresh placeholders and substitutions.
The construction therefore terminates. Every final predicate occurrence
comes from an input elementary label; the placeholder alphabets disappear.
No density, absence of endpoints, greatest lower bound or least upper
bound was needed for the fusion identity. The final property retains
its declared \(\mathsf{Com}\) scope, and no existential composition
from the next source Lemma~5 has been assumed.
\end{proof}
\end{annexitem}
